\paragraph{$^{$FOOTNOTE_NUMBER$}$The progenitor/descendant index}\label{footnote:$FOOTNOTE_NUMBER$} of a subhalo points to its progenitor/descendant in the SOAP catalogue of the previous/next snapshot. The progenitor index points to the subhalo that has the same HBT TrackId, and is $-1$ if that subhalo was not resolved in the previous snapshot. It can therefore only be used to move up the main progenitor branch, and provides no information about subhalo mergers.

The descendant index includes merger information. What happens to each subhalo between this snapshot and the next is given by \verb|DescendantType|, which is determined from the entry with the same TrackId in the HBT catalogue of the next snapshot:
\begin{itemize}
\item $0$ (survives): the subhalo is still resolved. The descendant is the subhalo with the same TrackId.
\item $1$ (disrupts): the subhalo is no longer self-bound. The descendant is the subhalo which contains the majority of its most bound tracer particles (HBT \verb|DescendantTrackId|).
\item $2$ (sinks): the core of the subhalo overlaps in phase-space with another subhalo (HBT \verb|SinkTrackId| $\neq -1$). The descendant is given by HBT \verb|DescendantTrackId|.
\item $3$ (disrupts with no descendant): the subhalo disrupts or sinks, but the majority of its most bound tracer particles are not bound to any subhalo (HBT \verb|DescendantTrackId| $= -1$). The descendant index is $-1$.
\item $-1$: the HBT catalogue for the next snapshot was not available (e.g. this is the final snapshot). The descendant index is $-1$.
\end{itemize}
